% Reviewer-style feedback for paper/user.tex.
% Written 2026-07-21. This document was not compiled to PDF.

\documentclass[11pt]{article}
\usepackage[margin=1in]{geometry}
\usepackage{kotex}
\usepackage[T1]{fontenc}
\usepackage{newtxtext,newtxmath}
\usepackage{booktabs}
\usepackage{longtable}
\usepackage{enumitem}
\usepackage{xcolor}
\usepackage[hidelinks]{hyperref}
\usepackage{url}

\definecolor{fatalred}{RGB}{150,20,20}
\definecolor{minorblue}{RGB}{20,70,130}
\newcommand{\fatal}{\textbf{\textcolor{fatalred}{[치명적]}}}
\newcommand{\minor}{\textbf{\textcolor{minorblue}{[사소]}}}
\newcommand{\fileline}[1]{\texttt{#1}}

\title{RelCompat3D 사용자 원고 검토}
\author{Reviewer-style transcript review}
\date{2026-07-21}

\begin{document}
\maketitle

\section{검토 범위와 총평}

검토 대상은 \fileline{paper/user.tex}에 조립한 사용자의 Abstract부터
Conclusion까지의 원고, Figure~1--3, Table~1--3이다. 실제 제출 원고인
\fileline{paper/aaai/} 아래 파일은 변경하지 않았다. PDF도 빌드하지 않았다.

현재 원고의 중심 논리는 대체로 정합하다.

\begin{enumerate}[leftmargin=2em]
    \item 고득점 relation score가 같은 ordered pair의 predicate--geometry
    compatibility를 직접 나타내지는 않는다는 failure를 정의한다.
    \item predictor score를 compatibility input에서 분리하고, pair identity와
    relation-preserving transformation을 보존하는 compatibility estimator를
    학습한다.
    \item proximity와 vertical-order만 family 내부에서 다시 정렬하고,
    support/contact는 source order로 유지한다.
    \item exact-match Recall과 verifier-derived Violation을 함께 보고하고,
    point/mesh audit와 controls로 construct dependence를 점검한다.
\end{enumerate}

따라서 task, method, main evidence의 큰 방향은 서로 맞는다. 다만 제출 전에 반드시
고쳐야 할 문제가 있다. 가장 큰 문제는 (i) Abstract의 사실관계 오류, (ii) SGFN
$K=5$ tie를 ``lower'' 또는 ``decreases''로 쓴 결과 과장, (iii) Figure caption과
실제 Figure의 불일치, (iv) Table~3 coverage의 모호성, (v) AAAI caption/figure
형식 위반이다. 독자가 이해하기 어려운 부분은 주로 Method의 첫 문단, counterfactual
설명, MLP 설명, family-aware re-ranking 절차, 그리고 Point/mesh audit에 집중되어
있다.

\subsection{심각도 기준}

\begin{itemize}[leftmargin=2em]
    \item \fatal{} 사실관계, claim, method scope, metric 해석, 또는 AAAI 형식 준수에
    직접 영향을 주어 제출 전 수정이 필요한 항목이다.
    \item \minor{} 핵심 결론을 바꾸지는 않지만 문법, collocation, 가독성, 용어
    통일성을 떨어뜨리는 항목이다.
\end{itemize}

\section{최우선 수정 목록}

\begin{longtable}{@{}p{0.10\textwidth}p{0.24\textwidth}p{0.58\textwidth}@{}}
\toprule
심각도 & 위치 & 문제와 조치 \\
\midrule
\endfirsthead
\toprule
심각도 & 위치 & 문제와 조치 \\
\midrule
\endhead
\fatal{} & Abstract & ``preserves ordered-pair measurements and the relation
score''는 실제 factor separation을 잘못 설명한다. ordered-pair \emph{identity}를
보존하고, $T/G$로 compatibility를 추정하며, $Z$는 re-ranking에서만 결합한다고
고쳐야 한다. \\
\fatal{} & Introduction/Conclusion & ``improve'' 및 ``lower ... at every
$K$''는 SGFN $K=5$ tie와 모순된다. ``improve or tie'', 또는 ``no worse point
estimates''로 고쳐야 한다. \\
\fatal{} & Results & ``SGFN is tied at $K=$''는 미완성 표기이며 $K=5$로
고쳐야 한다. 같은 문장의 ``Violation ... decreases for every predictor''도
``decreases or ties''가 정확하다. \\
\fatal{} & Figure~1 & 실제 그림은 strict elevation projection이 아니라
reconstructed point-cloud view이며, 6에서 425로의 이동은 Linear variant 결과다.
caption과 그림 내부의 ``Relcompat3D'' 표기를 수정해야 한다. \\
\fatal{} & Figure~2 & 그림 내부에는 transformation averaging이 나타나지 않는데
caption은 transformation-consistent compatibility를 말한다. 최소한 caption에서
averaging과 Linear rank example을 명시하고, primary re-ranking scope도 한 문장으로
한정해야 한다. \\
\fatal{} & Experiments & support/contact까지 포함한 scope를 ``can be re-ranked''로
설명한다. 실제로는 ``can be assessed''이며, re-ranking은 proximity/vertical-order에
한정된다. \\
\fatal{} & Table~3 & ``Ours''와 ``Cov.''가 모호하다. Ours를 Linear로 바꾸고,
coverage가 Linear의 measured/decidable coverage임을 명시해야 한다. Table~1의
Violation과 직접 비교할 수 없는 별도 construct라는 문장도 caption에 필요하다. \\
\fatal{} & AAAI format & Figure/Table caption의 bold lead-in, Figure~3의
\texttt{trim}/\texttt{clip}, 그리고 현재 root Figure PDF의 Identity-H fonts는
AAAI-27 규정과 충돌한다. \\
\fatal{} & Contribution~2 & transformation averaging 자체는 표준적인 group
averaging 성질이다. 독립 novelty처럼 쓰지 말고 source-score exclusion,
identity-preserving linked counterfactual supervision, exact transformation
consistency의 결합으로 claim해야 한다. \\
\bottomrule
\end{longtable}

\section{받은 피드백에 대한 판단}

\subsection{처음 읽는 독자에게 불친절한 부분}

이 피드백은 맞다. 특히 다음 다섯 부분에서 독자가 정의를 뒤늦게 조합해야 한다.

\begin{enumerate}[leftmargin=2em]
    \item Introduction에서 ``linked positive--counterfactual pairs''가 처음
    나오지만, counterfactual이 wrong pair인지 inverse predicate인지 설명되지 않는다.
    
    \item Problem Formulation의 첫 문단은 $r_i$, 두 identity tuple, $T_i$, $a_i$,
    $G_i$, score type을 한 번에 정의한다. 세 문단으로 분리하는 편이 낫다.

    \item ``family-aware re-ranking''이 왜 family label sequence를 보존하는지
    formal list 정의 전에는 직관적으로 설명되지 않는다. 작은 toy sequence 또는
    한 문장짜리 procedural description이 필요하다.

    \item ``predicate-linear skip connection''과 ``two-hidden-unit ReLU''는
    architecture에 익숙하지 않은 독자에게 역할이 보이지 않는다. predicate indicator가
    output으로 직접 가는 linear path라는 설명을 붙여야 한다.

    \item Point/mesh audit에서 ``robust vertex distances''와 ``area-weighted
    triangle samples''가 왜 OBB와 다른 construct인지 한 문장으로 풀어야 한다.
\end{enumerate}

원고를 장황하게 만들 필요는 없다. 각 지점에 한 문장씩만 보충하면 충분하다. 예를
들어 counterfactual은 다음처럼 처음 정의할 수 있다.

\begin{quote}
Linked counterfactuals pair each training positive with a deliberately
incompatible predicate or ordered pair constructed from the same training
context.
\end{quote}

family-aware re-ranking은 다음 한 문장으로 먼저 풀 수 있다.

\begin{quote}
At each source position, the method selects the next candidate from the same
relation family, but proximity and vertical-order candidates are reordered by
the combined score.
\end{quote}

\subsection{먼저 공개하지 않아도 되는 약점}

``약점을 숨긴다''와 ``main text에서 불필요한 future-work 목록을 줄인다''는 구분해야
한다. 다음 세 boundary는 결과 해석에 직접 영향을 주므로 숨기면 안 된다.

\begin{itemize}[leftmargin=2em]
    \item 세 predictor가 하나의 3DSSG/3RScan validation target에서 평가된다는 점.
    \item primary Violation이 verifier-derived이며 point/mesh audit도 독립적인
    physical-validity ground truth가 아니라는 점.
    \item support/contact는 평가하지만 primary method가 re-rank하지 않는다는 점.
\end{itemize}

반면 다음 내용은 main text에서 줄이거나 supplement로 옮겨도 된다.

\begin{itemize}[leftmargin=2em]
    \item ReplicaSSG/FROSS의 negative transfer를 구체적으로 먼저 말하는 문장. Main
    Discussion에는 single-target scope만 남기고 stress-test 결과는 supplement에서
    설명할 수 있다.
    \item ``Broader claims require ...''로 시작하는 reference labels/contact/pose/
    datasets의 긴 future-work 목록. 현재 audit이 independent ground truth가 아니라는
    한 문장만 남기면 충분하다.
    \item Metrics 문단의 pessimistic Violation, uncertainty rate, 두 종류 coverage
    수식 전체. Primary denominator는 본문에 유지하고 나머지는 supplement pointer로
    압축할 수 있다.
    \item Support/contact counterfactual 생성 세부 규칙. Primary variants가 이
    family를 re-rank하지 않으므로, all-family comparison에만 쓰인다는 설명 없이
    main Method에 두면 혼란스럽다.
    \item Re-ranking의 summed-utility maximization과 exchange-proof 설명. 보존
    property만 본문에 두고 formal claim은 supplement로 보낼 수 있다.
\end{itemize}

\subsection{Table~3에서 CI와 coverage를 뺄 것인가}

\textbf{판단: 95\% CI는 유지하는 것이 좋고, coverage도 원칙적으로 필요하다.}

95\% CI는 임의의 ``range''가 아니라 paired change의 sampling uncertainty를 보여준다.
NeurIPS 2024의 \emph{Verifiably Robust Conformal Prediction}은 main table에서
coverage와 95\% confidence interval을 함께 보고한다. CVPR 계열 논문도 metric
mean과 95\% CI lower/upper를 table에 두는 사례가 있다. 따라서 top-tier 관행을
이유로 CI를 제거할 근거는 없다.

\begin{itemize}[leftmargin=2em]
    \item NeurIPS example:
    \url{https://proceedings.neurips.cc/paper_files/paper/2024/file/0814a342597d65e0832fc7ec9b42c317-Paper-Conference.pdf}
    \item CVPR example:
    \url{https://openaccess.thecvf.com/content/CVPR2026F/supplemental/Li_Counterfactual_Segmentation_Reasoning_CVPRF_2026_supplemental.pdf}
\end{itemize}

coverage는 이 논문에서 더 직접적인 역할이 있다. Point/mesh disagreement가 uncertain이
되고, uncertain은 Violation numerator에는 들어가지 않기 때문에, Violation 감소가
decidable coverage 감소로 생긴 것이 아닌지 확인해야 한다. Selective-prediction
연구가 risk와 coverage를 함께 보고하는 이유와 같다.

\begin{itemize}[leftmargin=2em]
    \item NeurIPS selective-prediction example:
    \url{https://proceedings.neurips.cc/paper_files/paper/2023/file/a8526465a91166fbb90aaa8452b21eda-Paper-Conference.pdf}
\end{itemize}

다만 현재 표의 coverage는 불친절하다. \fileline{surface_audit/summary.json}을
확인하면 95.5/71.0, 98.2/83.7, 95.8/79.9는 Linear ranking의 measured/decidable
coverage다. Source coverage와의 차이는 measured coverage에서 0.5 percentage
point 이내, decidable coverage에서 1.3 points 이내다. 이를 caption 또는 본문에서
명시하면 coverage-collapse 공격을 더 직접적으로 막을 수 있다.

권장 header는 다음과 같다.

\begin{verbatim}
Predictor & Source V & Linear V & Delta V [95% CI] & Linear Cov. (M/D)
\end{verbatim}

지면 때문에 반드시 한 열을 줄여야 한다면, CI보다 coverage를 supplement로 옮기되
본문에 Source/Linear coverage가 크게 변하지 않았다는 수치를 한 문장으로 남기는 것이
낫다. 아무 설명 없이 coverage만 삭제하는 것은 권하지 않는다.

\subsection{Figure caption의 \texttt{\textbackslash texttt}}

\texttt{\textbackslash texttt} 자체가 학술 caption에서 금지된 명령은 아니다.
공식 ACL style example도 Figure caption 안에서 package 이름
\texttt{mwe}를 monospace로 표시한다.

\begin{itemize}[leftmargin=2em]
    \item Official ACL template:
    \url{https://github.com/acl-org/acl-style-files/blob/master/acl_latex.tex}
\end{itemize}

하지만 그 용도는 code, filename, command, literal token과 같이 monospace가 의미를
가질 때다. \texttt{heater close by trash can}은 code가 아니라 자연어 relation
triplet이므로 현재 caption에서는 이질적이다. 더구나 AAAI-27 Author Kit은 figure
caption을 10-point roman으로 두라고 명시한다
(\fileline{paper/aaai/official/AnonymousSubmission2027.tex:581}). 따라서 이
caption에서는 다음 중 하나를 권장한다.

\begin{verbatim}
Open3DSG ranks ``heater close by trash can'' at 19, ...
\end{verbatim}

또는 predicate만 italicize한다.

\begin{verbatim}
Open3DSG ranks the relation heater -- close by -- trash can at 19, ...
\end{verbatim}

\subsection{Contribution~2 재판단}

현재 Contribution~2는 그대로 두면 독립 contribution으로 약하게 읽힐 수 있다.

\begin{verbatim}
We learn predicate--geometry compatibility from the corresponding ordered
subject--object pair without using the predictor score, and enforce invariance
under applicable relation-preserving endpoint and predicate transformations.
\end{verbatim}

문제는 transformation averaging의 invariance 자체가 표준적인 group averaging
성질이라는 점이다. ``averaging을 사용했다''만으로는 novelty가 아니다. 그러나
다음 세 요소를 failure cause에서 유도된 하나의 method design으로 묶으면 contribution
수준이 된다.

\begin{enumerate}[leftmargin=2em]
    \item source relation score와 predictor identity를 compatibility input에서
    배제한다.
    \item candidate의 ordered-pair identity를 보존한 linked positive--counterfactual
    supervision을 사용한다.
    \item proximity/vertical-order의 의미 보존 변환에 대해 exact consistency를
    inference에서 보장한다.
\end{enumerate}

권장 Contribution~2는 다음과 같다.

\begin{quote}
We introduce source-score-excluded compatibility learning for fixed relation
candidates, using linked positive--counterfactual supervision and
relation-preserving transformation averaging to assign the same compatibility
to equivalent endpoint/predicate representations.
\end{quote}

이 표현은 transformation averaging만을 새롭다고 주장하지 않으면서도 factor
separation과 training supervision을 하나의 재현 가능한 method contribution으로
만든다. Contribution~3은 family composition과 unsupported family를 보존하는 constrained
re-ranking으로 별도 유지할 수 있다.

\section{섹션별 상세 검토}

\subsection{Title}

\paragraph{\minor{} ``Semantic Confidence''의 정밀도.}

원문:
\begin{verbatim}
Beyond Semantic Confidence: Relation-Consistent Geometric Re-ranking for
3D Scene Graphs
\end{verbatim}

제목의 후반부는 method 및 scope와 잘 맞는다. 다만 $Z$는 calibrated confidence가
아니라 sigmoid relation score 또는 cosine similarity이므로 ``Semantic Confidence''는
조금 넓다. 현재 제목은 rhetorical framing으로 유지 가능하지만, 최대한 엄밀하게 하려면
``Beyond Relation Scores''가 대안이다. 제목 변경은 필수는 아니다.

\subsection{Abstract}

\paragraph{\fatal{} 보존 대상과 factor separation의 사실관계 오류.}

원문:
\begin{verbatim}
RelCompat3D, a re-ranking framework that preserves ordered-pair measurements
and the relation score produced by the predictor.
\end{verbatim}

Method가 보존하는 핵심은 ordered-pair identity다. Measurements와 score를 하나로
``preserves''한다고 쓰면 둘을 model input으로 함께 유지하는 것처럼 읽힌다. 실제로
compatibility는 $T,G$로 계산하고 $Z$는 마지막 ranking에서만 결합한다.

수정 제안:
\begin{quote}
RelCompat3D preserves ordered-pair identity and estimates compatibility from
predicate semantics and predicate-independent ordered-pair measurements without
using predictor identity or score.
\end{quote}

\paragraph{\minor{} evaluation target 표현.}

원문:
\begin{verbatim}
on a shared 3DSSG
\end{verbatim}

3DSSG가 dataset, ontology, validation target 중 무엇인지 불명확하다.

수정 제안:
\begin{quote}
on one shared 3DSSG validation target
\end{quote}

\paragraph{\minor{} predictor--$K$ values collocation.}

원문:
\begin{verbatim}
Across all reported predictor--K values
\end{verbatim}

predictor는 value가 아니므로 ``settings''가 자연스럽다.

수정 제안:
\begin{quote}
Across all reported predictor--$K$ settings
\end{quote}

\paragraph{\minor{} alternative measure의 수 불일치.}

원문:
\begin{verbatim}
under an alternative geometric measure
\end{verbatim}

point와 mesh 두 representation을 결합한 audit이므로 ``using alternative geometric
measurements''가 더 정확하다.

\subsection{Introduction}

\paragraph{\fatal{} 빈 citation과 caption 형식.}

원문:
\begin{verbatim}
Open3DSG (source)~\cite{} predicts ...
\end{verbatim}

빈 citation은 제거하거나 실제 Open3DSG key를 넣어야 한다. 또한 모든 figure/table
caption의 bold lead-in은 AAAI-27의 10-point roman caption 규정과 맞지 않는다.
\texttt{\textbackslash textbf\{RelCompat3D.\}}를 삭제한다.

\paragraph{\fatal{} Figure~1의 view와 estimator가 부정확하다.}

원문:
\begin{verbatim}
although the elevation view places the desk below the ceiling.
...
RelCompat3D demotes the inconsistent relation to rank 425
\end{verbatim}

현재 Figure~1은 axis-aligned elevation plot이 아니라 reconstructed point-cloud view다.
또한 rank 425는 RelCompat3D-Linear 결과다.

수정 제안:
\begin{quote}
Open3DSG ranks ``desk higher than ceiling'' at 6, although the reconstructed
point-cloud view places the desk below the ceiling. RelCompat3D-Linear demotes
the same candidate to rank 425, outside the top 50.
\end{quote}

\paragraph{\fatal{} ``plausible'' motivation과 teaser example의 불일치.}

원문:
\begin{verbatim}
relation predictors rank plausible labels ... Figure 1 shows this failure.
\end{verbatim}

``desk higher than ceiling''은 category level에서도 자연스럽게 plausible한 relation은
아니다. Figure~1은 high-scoring geometric contradiction을 잘 보여주지만 category-level
plausibility를 직접 보여주지는 않는다. 첫 문장을 ``high-scoring labels''로 바꾸면
teaser와 claim이 정렬된다. Category-plausible failure는 Figure~2의 heater--trash-can
case로 설명할 수 있다.

\paragraph{\minor{} 핵심 method paragraph의 정보 밀도.}

원문:
\begin{verbatim}
Both estimators are trained with linked positive--counterfactual pairs ...
We then average compatibility over known relation-preserving transformations ...
\end{verbatim}

처음 읽는 독자는 counterfactual과 transformation이 각각 negative construction과
equivalence augmentation이라는 점을 구분하기 어렵다. 두 문장 사이에 다음 문장을
넣는 것이 좋다.

\begin{quote}
Counterfactuals create incompatible training pairs, whereas the transformations
encode alternative representations of the same relation semantics.
\end{quote}

\paragraph{\fatal{} all-$K$ result 문장의 overclaim.}

원문:
\begin{verbatim}
Across all reported K values, both variants improve the Recall--Violation
point-estimate trade-off over the source ranking
\end{verbatim}

SGFN $K=5$는 두 metric 모두 Source와 tie다. ``improve''를 strict improvement로
읽을 수 있으므로 다음이 정확하다.

\begin{quote}
Across all reported $K$ values, both variants yield no worse Recall--Violation
point estimates than the source rankings, with the largest changes on Open3DSG.
\end{quote}

\paragraph{\minor{} Contribution~1의 ``joint''가 combined metric처럼 읽힘.}

원문:
\begin{verbatim}
joint Recall@K--Violation@K reporting
\end{verbatim}

두 metric을 하나로 합친 objective가 아니므로 ``reporting Recall@$K$ and
Violation@$K$ together'' 또는 ``paired reporting''가 더 직접적이다.

\paragraph{\fatal{} Contribution~2 scope.}

원문:
\begin{verbatim}
enforce invariance under applicable relation-preserving endpoint and predicate
transformations
\end{verbatim}

``applicable''만으로는 support/contact까지 변환하는 것으로 읽힐 수 있다.
``the defined proximity and vertical-order transformations''로 범위를 명시한다.
Contribution-level 판단은 앞 절의 composite formulation을 따른다.

\subsection{Related Work}

\paragraph{\minor{} typo와 잘못된 collocation.}

원문:
\begin{verbatim}
absentation decisions
\end{verbatim}

``abstention decisions''가 맞다. 더 자연스럽게는 ``selective abstention''이다.

\paragraph{\minor{} Recall@$K$ 정의가 부정확함.}

원문:
\begin{verbatim}
Recall@K measures retrieval rank rather than compatibility ...
\end{verbatim}

Recall@$K$는 rank 자체가 아니라 cutoff에서의 exact-label retrieval을 측정한다.

수정 제안:
\begin{quote}
Recall@$K$ measures exact-label retrieval at a ranking cutoff rather than
compatibility with reconstructed ordered-pair geometry.
\end{quote}

\paragraph{\minor{} 2D prior-work 문장.}

원문:
\begin{verbatim}
For fixed 2D scene graph generation predictions, Neau et al. ...
\end{verbatim}

수식어 연결이 어색하다.

수정 제안:
\begin{quote}
For fixed predictions from 2D scene graph generators, Neau et al. ...
\end{quote}

\paragraph{잘 된 점.}

각 prose paragraph는 related methods와 RelCompat3D의 공통점 또는 차이로 끝난다.
3D generation, geometry-aware evidence, transformation methods, calibration/uncertainty와의
novelty boundary도 분리되어 있다. Related Work의 구조는 유지하는 것이 좋다.

\subsection{Method}

\paragraph{\minor{} Problem Formulation 첫 문단이 지나치게 조밀함.}

원문은 $r_i$, score contracts, context, 두 identity, $T_i$, $a_i$, $G_i$를 한
문단에서 정의한다. 다음 세 단위로 나누는 것이 좋다.

\begin{enumerate}[leftmargin=2em]
    \item candidate tuple, native source score, context;
    \item ordered-pair identity와 exact relation identity;
    \item $T_i$, $a_i$, $G_i$의 역할.
\end{enumerate}

\paragraph{\minor{} family label 설명의 문법과 의미.}

원문:
\begin{verbatim}
The family label a_i always selects the transformation set and re-ranking
relation families.
\end{verbatim}

하나의 $a_i$가 ``relation families''를 선택한다는 표현은 어색하다.

수정 제안:
\begin{quote}
The family label $a_i$ selects the applicable transformation set and determines
whether the candidate's family is re-ranked.
\end{quote}

\paragraph{\fatal{} support/contact training head의 목적이 빠짐.}

원문:
\begin{verbatim}
Support/contact negatives replace one endpoint under separation margins.
\end{verbatim}

primary variants는 support/contact를 re-rank하지 않는데 이 family의 negative
construction이 갑자기 등장한다. 이 detail을 supplement로 옮기거나, all-family
comparison의 linear head에만 쓰인다고 명시해야 한다.

\paragraph{\fatal{} exact construction rule pointer가 없음.}

원문:
\begin{verbatim}
Evaluation rows, predictor scores, and verifier-status labels are not used to
construct training examples.
\end{verbatim}

boundary는 좋지만 threshold, cap, balancing이 어디에 있는지 rendered text에서 알 수
없다. 다음 한 문장을 유지해야 한다.

\begin{quote}
The supplement specifies the family-specific thresholds, negative cap, and
balancing rules.
\end{quote}

이는 불필요한 약점 공개가 아니라 재현성 pointer다.

\paragraph{\minor{} nonlinear architecture 설명의 난이도.}

원문:
\begin{verbatim}
one shared two-hidden-unit ReLU network
...
With a predicate-linear skip connection
\end{verbatim}

``two-hidden-unit''는 hidden layer가 두 개인지 unit이 두 개인지 잠깐 모호하다.

수정 제안:
\begin{quote}
one shared single-hidden-layer ReLU network with two hidden units
\end{quote}

그리고 skip connection은 ``a direct linear path from predicate indicators to
the output logit''로 한 번 풀어 쓴다.

\paragraph{\minor{} outputs collocation.}

원문:
\begin{verbatim}
Transformation averaging therefore outputs ...
\end{verbatim}

수학적 relation에는 ``yields'' 또는 ``guarantees''가 자연스럽다.

\paragraph{\fatal{} loss가 실제 ordering을 보장하는 것처럼 씀.}

원문:
\begin{verbatim}
the pairwise objective still orders each linked positive above its
counterfactual
\end{verbatim}

loss는 ordering을 장려하지만 모든 sample에 대해 보장하지 않는다.

수정 제안:
\begin{quote}
the pairwise objective encourages each linked positive to score above its
counterfactual
\end{quote}

\paragraph{\fatal{} log-score 설명의 전제가 불명확함.}

원문:
\begin{verbatim}
For positive factors ... log u_i = log Z_i + log C_i ...
\end{verbatim}

Open3DSG의 native score는 cosine similarity이므로 score positivity가 일반 정의에서
자동으로 보장되지 않는다. 이 identity는 method에 필요하지 않으므로 삭제하는 것이
가장 안전하다. 유지하려면 $Z_i>0$ 및 $C_i^{\rm tr,q}>0$ 조건과 실제 observed score
contract를 명시해야 한다.

\paragraph{\minor{} optimization property가 main flow를 방해함.}

원문:
\begin{verbatim}
maximizes their sum of ranking scores subject to the preserved family counts
and support/contact subsequence
\end{verbatim}

construction property는 맞을 수 있지만 핵심 이해에는 필요하지 않다. 본문은 family
sequence와 support/contact subsequence 보존만 설명하고 exchange proof와 utility
maximization은 supplement로 옮기는 것이 읽기 쉽다.

\paragraph{\minor{} source score 명칭 통일.}

원고는 ``predictor score'', ``source score'', ``source relation score'', ``relation
score produced by the predictor''를 혼용한다. $Z$의 통일 표현은 ``source relation
score''가 가장 정확하다. probability 또는 calibrated confidence로 부르지 않는다.

\subsection{Experimental Setup}

\paragraph{\fatal{} evaluation scope와 re-ranking scope의 혼동.}

원문:
\begin{verbatim}
restricted to relations whose consistency can be re-ranked from reconstructed
ordered-pair geometry
\end{verbatim}

support/contact는 평가 scope에는 있지만 re-ranking scope에는 없다.

수정 제안:
\begin{quote}
restricted to relation families whose consistency can be assessed from
reconstructed ordered-pair geometry
\end{quote}

\paragraph{\minor{} family 이름 통일.}

원문:
\begin{verbatim}
relative vertical families
\end{verbatim}

Method와 맞춰 ``vertical-order family''로 통일한다.

\paragraph{\minor{} pooled model의 위치가 불명확함.}

원문:
\begin{verbatim}
a pooled linear model tests whether family-specific fitting is necessary
\end{verbatim}

Table~1에는 pooled model이 없다. ``reported in the supplement''를 바로 붙여야 독자가
누락된 baseline으로 오해하지 않는다.

\paragraph{\minor{} Metrics 문단의 과밀함.}

Primary Recall/Violation 수식과 uncertain denominator는 본문에 필요하다. 반면
uncertainty rate, decidable-only violation, status coverage, decidable coverage,
pessimistic variant를 한 문장에 모두 나열하면 main setup이 audit checklist처럼 읽힌다.

수정 제안:
\begin{quote}
The supplement reports uncertainty, coverage, decidable-only Violation, and a
pessimistic variant that counts uncertain candidates as violations.
\end{quote}

\paragraph{bootstrap 표현에 대한 판단.}

원문:
\begin{verbatim}
Paired 95% confidence intervals use 1,000 bootstrap resamples of the 157 scans,
keeping all contexts from each scan together ...
\end{verbatim}

이 문장은 필요하고 자연스럽다. ``bootstrap''은 top-tier empirical papers에서 흔히
쓰는 표준 통계 용어다. 같은 scan의 context가 독립 sample이 아니므로 scan-level
resampling을 밝히는 것이 오히려 rigor를 높인다. ``cluster bootstrap''이라는 말을
반드시 쓸 필요는 없고, 현재처럼 procedural definition을 쓰면 충분하다.

\subsection{Results}

\paragraph{\fatal{} tie를 decrease로 서술함.}

원문:
\begin{verbatim}
the Violation point estimate decreases for every predictor. SGFN is tied at
K= and changes only marginally at K=10
\end{verbatim}

SGFN $K=5$는 Source, Linear, MLP가 Recall 31.17, Violation 2.37로 모두 tie다.

수정 제안:
\begin{quote}
Across $K\in\{5,10,20,50,100\}$, neither variant has a lower Recall point
estimate or a higher Violation point estimate than Source. SGFN is tied at
$K=5$ and changes only marginally at $K=10$.
\end{quote}

\paragraph{\minor{} ``reported ranges'' collocation.}

원문:
\begin{verbatim}
the largest gains across the reported ranges
\end{verbatim}

``reported $K$ values''가 정확하다.

\paragraph{\minor{} ``better''의 정의.}

원문:
\begin{verbatim}
yields a better K=50 point estimate Recall-Violation trade-off
\end{verbatim}

두 metric이 모두 weakly 개선된다는 뜻을 바로 밝히면 좋다.

수정 제안:
\begin{quote}
at $K=50$, each variant has a Recall point estimate no lower and a Violation
point estimate lower than Source for all three predictors
\end{quote}

\paragraph{잘 된 점.}

Linear와 MLP 중 하나를 universal winner로 만들지 않고 Open3DSG의 Recall--Violation
operating-point 차이를 설명한 점은 엄밀하다. RankAvg/RRF의 low-$K$ Recall loss와
all-family product의 support/contact selection change도 comparator 역할과 맞게
해석했다.

\subsection{Discussion and Limitations}

\paragraph{\minor{} main text에서 줄일 수 있는 stress-test 문장.}

원문:
\begin{verbatim}
The supplementary ReplicaSSG evaluation using FROSS predictions shows that
performance changes under ontology and geometry shifts.
\end{verbatim}

정직하지만 main claim을 방어하는 데 필수는 아니다. single-target scope를 바로
말한 뒤 이 문장은 supplement에 두는 편이 submission narrative에는 더 유리하다.
삭제하더라도 cross-dataset generalization을 주장해서는 안 된다.

\paragraph{\fatal{} 반드시 남겨야 하는 construct boundary.}

원문:
\begin{verbatim}
The point- and mesh-based audit still uses the same reconstructed instance
geometry and relation ontology, so it does not provide an independent reference
for geometric validity.
\end{verbatim}

이 문장은 핵심 construct validity boundary이므로 유지해야 한다. 숨기면 reviewer가
audit을 independent ground truth로 오해하거나 circularity를 지적했을 때 더 큰 문제가
된다. 다만 다음 future-work 문장은 줄일 수 있다.

\begin{verbatim}
Broader claims require independently defined reference labels, richer contact
and pose evidence, and evaluation on additional datasets.
\end{verbatim}

권장 압축:
\begin{quote}
Because the point- and mesh-based audit uses the same reconstructed scenes and
ontology, it is an alternative geometric measurement rather than independent
validity ground truth.
\end{quote}

\paragraph{\minor{} outside-scope 목록 압축.}

support/contact는 main boundary이므로 유지한다. Reference-frame relations와 complete
graph generation은 한 문장으로 줄일 수 있다.

\begin{quote}
RelCompat3D assumes known instances and reconstructed ordered-pair geometry; it
does not generate complete graphs or correct contact-dependent relations.
\end{quote}

\subsection{Conclusion}

\paragraph{\fatal{} strict lower claim이 Table~1과 모순됨.}

원문:
\begin{verbatim}
both variants yield lower verifier-derived Violation point estimates while
preserving or improving Recall point estimates at every reported K
\end{verbatim}

SGFN $K=5$는 Violation도 tie다.

수정 제안:
\begin{quote}
Across three predictors on one shared 3DSSG validation target, both variants
yield lower or tied verifier-derived Violation point estimates while preserving
or improving Recall point estimates at every reported $K$.
\end{quote}

\section{Figure 검토}

\subsection{Figure~1: failure case}

\paragraph{내용 적합성.}

Figure~1은 failure case의 목적을 대체로 잘 수행한다. 같은 candidate가 Source rank
6에서 Linear rank 425로 이동하고, desk가 ceiling 아래에 있는 reconstructed point-cloud
evidence가 함께 보인다. 따라서 ``high source score does not ensure same-pair geometric
compatibility''라는 task motivation은 전달된다.

다만 다음 수정이 필요하다.

\begin{itemize}[leftmargin=2em]
    \item 그림 내부 ``Relcompat3D''를 ``RelCompat3D-Linear''로 수정한다.
    \item caption의 ``elevation view''를 ``reconstructed point-cloud view''로
    수정한다.
    \item rank change가 Linear 결과임을 caption에 명시한다.
    \item 이 사례는 category-plausible error라기보다 obvious vertical contradiction이다.
    Introduction의 ``plausible'' failure와 동일 사례라고 단정하지 않는다.
\end{itemize}

\subsection{Figure~2: method overview}

현재 Figure~2가 보여주는 핵심 흐름은 맞다.

\begin{enumerate}[leftmargin=2em]
    \item predicate semantics와 pair measurements만 compatibility block으로 들어간다.
    \item source relation score는 compatibility를 우회한다.
    \item 두 signal은 within-family score에서 결합된다.
    \item family-aware re-ranking 후 rank가 19에서 178로 이동한다.
\end{enumerate}

따라서 Linear/MLP architecture, BCE, margin loss를 그림 안에 추가할 필요는 없다.
그림을 과밀하게 만들지 않으려는 판단은 타당하다. 그러나 다음 두 정보는 caption에
필요하다.

\begin{itemize}[leftmargin=2em]
    \item compatibility는 relation-preserving transformations에 대해 averaging된
    score라는 점;
    \item illustrated 19--178 rank change는 RelCompat3D-Linear 결과라는 점.
\end{itemize}

가능하면 그림 내부 ``T/G Compatibility''는 ``Predicate--geometry compatibility''로
바꾼다. Slash 표기는 $T$와 $G$의 단순 ratio처럼 보일 수 있다. Primary scope를 그림에
더 넣기 어렵다면 caption 끝에 다음 한 문장을 추가한다.

\begin{quote}
The method re-ranks proximity and vertical-order candidates; support/contact
candidates retain source order. The illustrated 19-to-178 change is produced by
RelCompat3D-Linear.
\end{quote}

\subsection{Figure~3: quantitative trajectory}

Figure~3의 plot type은 적합하다. Table~1이 exact values를, Figure~3이 right/down
movement와 $K$에 따른 trajectory를 담당하므로 bar chart보다 현재 line trajectory가
더 많은 정보를 보존한다. 세 predictor의 scale이 달라 axis range가 다른 점도 caption에
명시되어 있다.

수정할 부분은 다음과 같다.

\begin{itemize}[leftmargin=2em]
    \item legend의 ``Ours (Linear)/(MLP)''를 가능하면 Table과 같은
    ``RelCompat3D-Linear/MLP''로 통일한다.
    \item ``same budget order''를 ``same increasing-$K$ order''로 바꾼다.
    \item 현재 45 coordinates는 active
    \fileline{evaluation/routed_comparators/metrics.csv}의 수치와 일치한다.
\end{itemize}

\subsection{Figure source의 AAAI 형식 문제}

\fatal{} \fileline{paper/user.tex}이 직접 참조하는
\fileline{paper/Figure1.pdf}, \fileline{Figure2.pdf}, \fileline{Figure3.pdf}를
\texttt{pdffonts}로 확인하면 모두 CID TrueType/Identity-H font를 포함한다. AAAI-27
Author Kit은 CID/Identity-H를 제거하거나 outline으로 변환하도록 요구한다
(\fileline{AnonymousSubmission2027.tex:232}). 또한 Figure~3 inclusion은 다음을
사용한다.

\begin{verbatim}
trim={0bp 43bp 0bp 45bp}, clip
\end{verbatim}

Author Kit은 LaTeX의 trim/clip에 의존한 crop을 금지하고 figure file 자체를 외부에서
crop하라고 명시한다
(\fileline{AnonymousSubmission2027.tex:575--579}). 따라서 최종 제출에서는
outlined-font PDF를 만들고 crop을 asset에 적용한 뒤 단순
\texttt{\textbackslash includegraphics[width=...]}로 삽입해야 한다.

\section{Table 및 수치 검토}

\subsection{Table~1}

\paragraph{수치.}

Table~1의 Source, Linear, MLP, RankAvg, RRF values는 active routed-comparator
artifact와 일치한다. Abstract의 ``no lower Recall/no higher Violation point
estimates''도 모든 30 proposed-variant predictor--$K$ pair에 대해 point-estimate
수준에서 맞다. 단, tie를 strict decrease로 바꾸면 안 된다.

\paragraph{\minor{} caption 문법.}

원문:
\begin{verbatim}
Source denoted each predictor's original ranking.
\end{verbatim}

``Source denotes''가 맞다. ``All entries are percentages''도 넣어야 한다. Caption의
bold lead-in은 제거한다.

권장 caption:
\begin{quote}
Shared 3DSSG validation results. Exact-match Recall (R, higher is better) and
verifier-derived Violation (V, lower is better) are percentages. Source denotes
each predictor's original ranking. Both RelCompat3D variants, RankAvg, and RRF
use the same family-aware ranking procedure; Product (all families) applies
compatibility to every evaluated family.
\end{quote}

\subsection{Table~2: Ablations}

\paragraph{수치.}

제시한 Linear control 값과 full MLP rows는 locked K=50/100 artifacts와 정합하다.
각 predictor block의 row count도 8로 맞다.

\paragraph{\fatal{} caption의 scope가 완전히 명시되지 않음.}

원문:
\begin{verbatim}
For compactness, control rows are shown for RelCompat3D-Linear.
\end{verbatim}

표 안의 MLP row가 full method인지 MLP control인지 caption만으로는 모호하다. 또한
모든 값이 percentage라는 설명이 빠졌다.

권장 caption:
\begin{quote}
Ablations and counterfactual controls at $K\in\{50,100\}$. Control rows use
RelCompat3D-Linear; the RelCompat3D-MLP rows report the full nonlinear variant
for comparison, with matched MLP controls in the supplement. All entries are
percentages; R is Recall (higher is better) and V is verifier-derived Violation
(lower is better). Every condition uses the same candidates and family-aware
ranking procedure as Table~1.
\end{quote}

\subsection{Table~3: Point/mesh audit}

\paragraph{수치.}

Source, Linear, paired change, and Linear measured/decidable coverage는 frozen
surface-audit artifact와 일치한다. Change는 percentage-point 차이다.

\paragraph{\fatal{} ``Ours''와 coverage provenance.}

원문:
\begin{verbatim}
Predictor & Source & Ours & Change (95% CI) & Cov.
...
Ours denotes RelCompat3D-Linear; ... Cov. is measured/decidable coverage.
\end{verbatim}

두 proposed variant가 있으므로 ``Ours''는 모호하다. ``Linear''로 바꾼다. Coverage가
Source인지 Linear인지도 header에서 드러나지 않는다. ``Linear Cov. (M/D)''로
바꾸고 caption에서 M/D를 푼다.

\paragraph{\fatal{} primary metric과 다른 construct라는 설명 누락.}

현재 body에는 설명이 있지만 caption만 읽으면 Table~1 Violation을 다른 방법으로
재측정한 것으로 오해할 수 있다. ``These values are not directly comparable to the
primary Violation in Table~1''를 caption에 넣는다.

권장 compact caption:
\begin{quote}
Point-/mesh-based Violation at $K=50$ (\%, lower is better). Satisfied or
violated status requires agreement between the two measurements; disagreement
is uncertain. $\Delta V$ is Linear minus Source in percentage points with a
paired 95\% interval from scan-level resampling; M/D is Linear
measured/decidable coverage. These values are not directly comparable to the
primary Violation in Table~1.
\end{quote}

이 caption이 너무 길면 agreement rule을 바로 앞 본문에 두되, alternative metric과
coverage provenance는 caption에 남긴다.

\section{표현 통일 계약}

다음 표현으로 통일하는 것이 좋다.

\begin{longtable}{@{}p{0.31\textwidth}p{0.61\textwidth}@{}}
\toprule
개념 & 통일 표현 \\
\midrule
generic field & 3D scene graph (고유명/제목이 아니면 lowercase) \\
$Z$ & source relation score \\
$G$ & predicate-independent ordered-pair measurements \\
$C$ & predicate--geometry compatibility \\
pair & 첫 등장 ordered subject--object pair, 이후 ordered pair \\
vertical family & vertical-order relations/candidates \\
metric & exact-match Recall@$K$; verifier-derived Violation@$K$ \\
ranking & family-aware re-ranking \\
support/contact & candidates retain source order \\
audit & point- and mesh-based consistency audit / alternative geometric measure \\
\bottomrule
\end{longtable}

다음 혼용은 정리한다.

\begin{itemize}[leftmargin=2em]
    \item predictor score / source score / source relation score
    \item predicate-geometry / predicate--geometry
    \item vertical / relative vertical / vertical order
    \item ordered pair geometry / ordered-pair geometry
    \item Recall@K / Recall@$K$
\end{itemize}

\section{권장 수정 순서}

\begin{enumerate}[leftmargin=2em]
    \item Abstract의 factor-separation 문장, Results의 SGFN tie, Conclusion의 strict
    lower claim을 먼저 수정한다.
    \item Figure~1/2 caption을 실제 view와 Linear rank change에 맞추고 caption의
    bold/monospace를 제거한다.
    \item Method의 support/contact negative 목적, counterfactual rule pointer,
    pairwise-loss 보장 표현, log-score 문장을 정리한다.
    \item Experiments의 assessed/re-ranked scope를 분리하고 Table~2/3 caption을
    self-contained하게 만든다.
    \item Discussion은 single-target와 non-independent-audit boundary만 남기고
    FROSS 상세 및 future-work 목록을 supplement로 축소한다.
    \item 최종 Figure PDF의 Identity-H fonts와 Figure~3 trim/clip을 제거한다.
\end{enumerate}

\section{최종 판단}

현재 원고는 ``geometry를 추가했다''는 단순 claim보다 강한 구조를 갖는다. Failure의
원인을 source relation score와 same-pair compatibility의 차이로 정의하고, 그 원인에
맞춰 score-excluded compatibility, identity preservation, transformation consistency,
family-constrained re-ranking을 설계했기 때문이다. Wrong-pair/shuffled-geometry,
wrong-predicate/swap, distance-only, compatibility-only controls도 이 논리를 검증한다.

그러나 Contribution~2를 transformation averaging 자체로 내세우면 novelty ceiling이
낮아진다. 반드시 source-score exclusion과 linked counterfactual supervision을 포함한
composite design으로 써야 한다. 또한 verifier-derived result를 strict physical-validity
claim으로 넓히지 않고, point/mesh audit을 alternative construct check로 한정해야 한다.
이 두 조건과 위의 사실/형식 오류를 수정하면, transcript는 reviewer가 task, method,
evidence, boundary를 일관되게 요약할 수 있는 수준에 도달한다.

\end{document}
